\documentclass{article}
\usepackage{amsmath,amssymb,amsthm}
\usepackage{algorithm}
\usepackage[noend]{algpseudocode}
\usepackage{fullpage}
\usepackage{hyperref}
\newtheorem{theorem}{Theorem}
\newtheorem{lemma}{Lemma}
\newtheorem{assumption}{Assumption}
\newcommand{\E}{\mathbb{E}}
\newcommand{\R}{\mathbb{R}}
\newcommand{\norm}[1]{\left\lVert#1\right\rVert}
\newcommand{\grad}{\nabla}
\begin{document}

\section*{Algorithm Name}
\textbf{Stochastic Subsampled Newton (SSN)}  

The method uses a stochastic (mini‑batch) estimate of the Hessian to construct a Newton‑type direction and updates the iterate with a stepsize $\alpha>0$.

\section*{Mathematical Formulation}
Let $f:\R^{d}\rightarrow\R$ be the empirical risk
\[
f(x)=\frac{1}{n}\sum_{i=1}^{n} f_i(x),\qquad 
f_i\in C^{2}.
\]
At iteration $t$ we draw two independent mini‑batches 
\[
\mathcal{B}^{(g)}_{t}\subset\{1,\dots,n\},\qquad 
\mathcal{B}^{(H)}_{t}\subset\{1,\dots,n\},
\]
with sizes $b_g$ and $b_H$, respectively.  

\noindent\textbf{Stochastic gradient}
\[
g_t\;:=\;\frac{1}{b_g}\sum_{i\in\mathcal{B}^{(g)}_{t}} \grad f_i(x_t).
\]

\noindent\textbf{Stochastic Hessian}
\[
\widehat H_t\;:=\;\frac{1}{b_H}\sum_{i\in\mathcal{B}^{(H)}_{t}} \grad^{2} f_i(x_t).
\]
We assume $\widehat H_t$ is invertible (or regularised, e.g. $\widehat H_t+\lambda I$).  

\noindent\textbf{Update rule}
\begin{equation}
\boxed{%
x_{t+1}=x_t-\alpha\,\widehat H_t^{-1} g_t
}\label{eq:update}
\end{equation}
where $\alpha\in(0,1]$ is a constant stepsize (or a diminishing schedule, see discussion).

\section*{Pseudocode}
\begin{algorithm}[H]
\caption{Stochastic Subsampled Newton (SSN)}\label{alg:ssn}
\begin{algorithmic}[1]
\State \textbf{Input:} stepsize $\alpha>0$, batch sizes $b_g,b_H$, max iter $T$
\State Initialize $x_0\in\R^{d}$
\For{$t=0,\dots,T-1$}
    \State Sample $\mathcal{B}^{(g)}_{t}\sim\text{Uniform}(\{1,\dots,n\})$ of size $b_g$
    \State Sample $\mathcal{B}^{(H)}_{t}\sim\text{Uniform}(\{1,\dots,n\})$ of size $b_H$
    \State $g_t\gets\frac{1}{b_g}\sum_{i\in\mathcal{B}^{(g)}_{t}}\grad f_i(x_t)$
    \State $\widehat H_t\gets\frac{1}{b_H}\sum_{i\in\mathcal{B}^{(H)}_{t}}\grad^{2} f_i(x_t)$
    \State \textbf{(Regularisation)} $\widehat H_t\gets\widehat H_t+\lambda I$ with $\lambda>0$ if needed
    \State $x_{t+1}\gets x_t-\alpha\,\widehat H_t^{-1} g_t$
\EndFor
\State \textbf{Output:} $x_T$ (or averaged iterate $\bar x_T=\frac1T\sum_{t=0}^{T-1}x_t$)
\end{algorithmic}
\end{algorithm}

\section*{Dry Run Example}
Consider the one‑dimensional quadratic 
\[
f(w)=(w-3)^{2}.
\]
Here $f'(w)=2(w-3)$ and $f''(w)=2$ (exact Hessian is constant).  
We use a deterministic (batch size $=1$) setting so that stochastic estimates coincide with the true quantities.  
Set $\alpha=0.1$, $w_0=0$ and run three iterations.

\begin{center}
\begin{tabular}{c|c|c|c}
$t$ & $g_t= f'(w_t)$ & $\widehat H_t = f''(w_t)$ & $w_{t+1}=w_t-\alpha\widehat H_t^{-1}g_t$ \\ \hline
0 & $2(0-3)=-6$ & $2$ & $0-0.1\cdot(1/2)(-6)=0+0.3=0.3$ \\[4pt]
1 & $2(0.3-3)=-5.4$ & $2$ & $0.3-0.1\cdot(1/2)(-5.4)=0.3+0.27=0.57$ \\[4pt]
2 & $2(0.57-3)=-4.86$ & $2$ & $0.57-0.1\cdot(1/2)(-4.86)=0.57+0.243=0.813$
\end{tabular}
\end{center}
Objective values:
\[
f(w_0)=9,\; f(w_1)=7.29,\; f(w_2)=5.44,\; f(w_3)=4.07.
\]
The iterates move toward the minimiser $w^{\star}=3$, confirming the Newton‑type scaling.

\newpage
\section*{Assumptions}
\begin{assumption}[Smoothness]\label{ass:smooth}
$f$ is $L$‑smooth, i.e.
\[
\forall x,y\in\R^{d}:\quad
f(y)\le f(x)+\langle\grad f(x),y-x\rangle+\frac{L}{2}\norm{y-x}^{2}.
\]
Equivalently $\norm{\grad f(x)-\grad f(y)}\le L\norm{x-y}$.
\end{assumption}

\begin{assumption}[Bounded stochastic gradient variance]\label{ass:gradvar}
There exists $\sigma_g^{2}>0$ such that for all $t$,
\[
\E\!\left[\norm{g_t-\grad f(x_t)}^{2}\mid x_t\right]\le\sigma_g^{2}.
\]
Moreover $\E[g_t\mid x_t]=\grad f(x_t)$ (unbiasedness).
\end{assumption}

\begin{assumption}[Unbiased Hessian estimator and bounded variance]\label{ass:hessian}
Define $H(x):=\frac{1}{n}\sum_{i=1}^{n}\grad^{2}f_i(x)$.  
There exists $\sigma_H^{2}>0$ with
\[
\E\!\left[\norm{\widehat H_t-H(x_t)}_{F}^{2}\mid x_t\right]\le\sigma_H^{2},
\qquad
\E[\widehat H_t\mid x_t]=H(x_t).
\]
\end{assumption}

\begin{assumption}[Positive definiteness of the true Hessian]\label{ass:posdef}
For all $x$, the true Hessian satisfies
\[
\mu I \preceq H(x) \preceq L I,
\]
with constants $0<\mu\le L$. Consequently $H(x)$ is invertible and $\norm{H(x)^{-1}}\le 1/\mu$.
\end{assumption}

\begin{assumption}[Stepsize]\label{ass:step}
The stepsize $\alpha$ is chosen as
\[
0<\alpha\le \min\!\left\{1,\frac{\mu}{L^{2}}\right\}.
\]
\end{assumption}

\section*{Main Convergence Theorem}
\begin{theorem}[Expected gradient norm decay]\label{thm:main}
Under Assumptions \ref{ass:smooth}–\ref{ass:step}, the iterates generated by \eqref{eq:update} satisfy for any $T\ge 1$
\[
\frac{1}{T}\sum_{t=0}^{T-1}\E\!\left[\norm{\grad f(x_t)}^{2}\right]
\;\le\;
\frac{2\bigl(f(x_0)-f^{\star}\bigr)}{\alpha\mu T}
\;+\;
\frac{2\alpha L\sigma_g^{2}}{\mu}
\;+\;
\frac{2\alpha^{2}L^{2}\sigma_H^{2}}{\mu^{2}},
\]
where $f^{\star}= \inf_{x}f(x)$. Consequently,
\[
\min_{0\le t<T}\E\!\left[\norm{\grad f(x_t)}^{2}\right]=O\!\left(\frac{1}{T}\right)
\quad\text{if }\alpha=O\!\left(\frac{1}{\sqrt{T}}\right).
\]
\end{theorem}

\section*{Theoretical Properties}
\begin{itemize}
\item \textbf{Stability.} The regularisation $\lambda I$ guarantees $\widehat H_t^{-1}$ is well‑defined; the bound $\norm{\widehat H_t^{-1}}\le 1/(\mu-\lambda)$ follows from Assumption \ref{ass:posdef}.
\item \textbf{Hyper‑parameter sensitivity.} The rate exhibits a linear dependence on $\alpha$ for the variance terms; choosing $\alpha\asymp T^{-1/2}$ balances the bias term $O(1/(\alpha T))$ with the variance terms $O(\alpha)$.
\item \textbf{Comparison to SGD.} For the same batch sizes, SSN improves the dependence on $L$ from $L^{2}$ (in SGD) to $L$ in the leading term because the Newton direction pre‑conditions the gradient.
\end{itemize}

\section*{Discussion}
The theorem shows that, despite non‑convexity, SSN achieves a \emph{stationarity} rate comparable to first‑order methods while exploiting curvature information to reduce the constant factor. If the Hessian variance $\sigma_H^{2}=0$ (exact Hessian or large $b_H$), the bound recovers the deterministic Newton rate $\mathcal{O}(1/(\alpha\mu T))$.

\newpage
\section*{Preliminaries (Definitions and Lemmas)}
\begin{lemma}[Descent Lemma for $L$‑smooth functions]\label{lem:descent}
For any $x$ and any vector $d$,
\[
f(x+d)\le f(x)+\langle\grad f(x),d\rangle+\frac{L}{2}\norm{d}^{2}.
\]
\end{lemma}
\begin{proof}
Directly from Assumption \ref{ass:smooth}. \hfill $\square$
\end{proof}

\begin{lemma}[Bound on inverse stochastic Hessian]\label{lem:hessian_inv}
Under Assumptions \ref{ass:posdef} and a regularisation $\lambda\in[0,\mu)$,
\[
\bigl\|\widehat H_t^{-1}\bigr\|\;\le\;\frac{1}{\mu-\lambda}\quad\text{a.s.}
\]
\end{lemma}
\begin{proof}
Because $H(x_t)\succeq\mu I$ and $\widehat H_t=H(x_t)+E_t$ with $E_t$ symmetric, adding $\lambda I$ yields $\widehat H_t+\lambda I\succeq (\mu-\lambda)I$, hence the spectral norm of the inverse is bounded by $(\mu-\lambda)^{-1}$. \hfill $\square$
\end{proof}

\section*{Main Proof (Step‑by‑Step Derivation)}
We prove Theorem \ref{thm:main}. Throughout, expectations are taken w.r.t. all randomness up to iteration $t$ unless otherwise noted.

\noindent\textbf{Step 1. Apply the descent lemma to the update \eqref{eq:update}.}  

Using Lemma \ref{lem:descent} with $d_t:=-\alpha\widehat H_t^{-1}g_t$,
\begin{align}
f(x_{t+1}) &\le f(x_t)+\big\langle\grad f(x_t),-\alpha\widehat H_t^{-1}g_t\big\rangle
+\frac{L}{2}\big\|\alpha\widehat H_t^{-1}g_t\big\|^{2}.
\tag{1}
\end{align}
[Step 1]

\noindent\textbf{Step 2. Insert unbiasedness of $g_t$.}  

Decompose $g_t=\grad f(x_t)+\xi_t$ where $\E[\xi_t\mid x_t]=0$ and $\E\!\left[\|\xi_t\|^{2}\mid x_t\right]\le\sigma_g^{2}$ (Assumption \ref{ass:gradvar}). Then
\[
\big\langle\grad f(x_t),\widehat H_t^{-1}g_t\big\rangle
=\big\langle\grad f(x_t),\widehat H_t^{-1}\grad f(x_t)\big\rangle
+\big\langle\grad f(x_t),\widehat H_t^{-1}\xi_t\big\rangle .
\tag{2}
\]
[Step 2]

\noindent\textbf{Step 3. Take conditional expectation of (1).}  

Condition on $x_t$ and use $\E[\xi_t\mid x_t]=0$:
\[
\E\!\bigl[f(x_{t+1})\mid x_t\bigr]\le
f(x_t)-\alpha\big\langle\grad f(x_t),\widehat H_t^{-1}\grad f(x_t)\big\rangle
+\frac{L\alpha^{2}}{2}\E\!\bigl[\|\widehat H_t^{-1}g_t\|^{2}\mid x_t\bigr].
\tag{3}
\]
[Step 3]

\noindent\textbf{Step 4. Lower bound the quadratic term using positive definiteness.}  

From Lemma \ref{lem:hessian_inv} and Assumption \ref{ass:posdef},
\[
\big\langle\grad f(x_t),\widehat H_t^{-1}\grad f(x_t)\big\rangle
\ge \frac{1}{L}\norm{\grad f(x_t)}^{2},
\tag{4}
\]
because $\widehat H_t^{-1}\succeq \frac{1}{L}I$ a.s. (since $\widehat H_t\preceq L I$).  
[Step 4]

\noindent\textbf{Step 5. Expand the squared norm in (3).}  

\[
\|\widehat H_t^{-1}g_t\|^{2}
=\|\widehat H_t^{-1}\grad f(x_t)\|^{2}
+2\langle\widehat H_t^{-1}\grad f(x_t),\widehat H_t^{-1}\xi_t\rangle
+\|\widehat H_t^{-1}\xi_t\|^{2}.
\tag{5}
\]
Taking conditional expectation and using $\E[\xi_t\mid x_t]=0$ kills the cross term:
\[
\E\!\bigl[\|\widehat H_t^{-1}g_t\|^{2}\mid x_t\bigr]
=\|\widehat H_t^{-1}\grad f(x_t)\|^{2}
+\E\!\bigl[\|\widehat H_t^{-1}\xi_t\|^{2}\mid x_t\bigr].
\tag{6}
\]
[Step 5]

\noindent\textbf{Step 6. Upper bound the two terms in (6).}  

Using Lemma \ref{lem:hessian_inv},
\[
\|\widehat H_t^{-1}\grad f(x_t)\|^{2}
\le \bigl\|\widehat H_t^{-1}\bigr\|^{2}\norm{\grad f(x_t)}^{2}
\le \frac{1}{\mu^{2}}\norm{\grad f(x_t)}^{2}.
\tag{7}
\]
Similarly,
\[
\E\!\bigl[\|\widehat H_t^{-1}\xi_t\|^{2}\mid x_t\bigr]
\le \bigl\|\widehat H_t^{-1}\bigr\|^{2}\E\!\bigl[\|\xi_t\|^{2}\mid x_t\bigr]
\le \frac{\sigma_g^{2}}{\mu^{2}}.
\tag{8}
\]
[Step 6]

\noindent\textbf{Step 7. Substitute (4), (7) and (8) into (3).}  

\[
\E\!\bigl[f(x_{t+1})\mid x_t\bigr]\le
f(x_t)-\frac{\alpha}{L}\norm{\grad f(x_t)}^{2}
+\frac{L\alpha^{2}}{2}\Bigl(\frac{1}{\mu^{2}}\norm{\grad f(x_t)}^{2}
+\frac{\sigma_g^{2}}{\mu^{2}}\Bigr).
\tag{9}
\]
[Step 7]

\noindent\textbf{Step 8. Collect the terms that contain $\|\grad f(x_t)\|^{2}$.}  

\[
\E\!\bigl[f(x_{t+1})\mid x_t\bigr]\le
f(x_t)-\underbrace{\Bigl(\frac{\alpha}{L}-\frac{L\alpha^{2}}{2\mu^{2}}\Bigr)}_{\displaystyle c_{\alpha}}
\norm{\grad f(x_t)}^{2}
+\frac{L\alpha^{2}\sigma_g^{2}}{2\mu^{2}}.
\tag{10}
\]
[Step 8]

\noindent\textbf{Step 9. Choose $\alpha$ to make $c_{\alpha}>0$.}  

Assumption \ref{ass:step} guarantees $\alpha\le \mu/L^{2}$, thus
\[
\frac{\alpha}{L}-\frac{L\alpha^{2}}{2\mu^{2}}
\ge \frac{\alpha}{2L}>0.
\tag{11}
\]
Hence we can set $c_{\alpha}\ge \frac{\alpha}{2L}$.  
[Step 9]

\noindent\textbf{Step 10. Take total expectation and sum over $t=0,\dots,T-1$.}  

Denote $\Delta_t:=\E[f(x_t)]-f^{\star}\ge0$. Using (10) and (11):
\[
\Delta_{t+1}\le\Delta_t-\frac{\alpha}{2L}\E\!\bigl[\norm{\grad f(x_t)}^{2}\bigr]
+\frac{L\alpha^{2}\sigma_g^{2}}{2\mu^{2}}.
\tag{12}
\]
Summing (12) from $t=0$ to $T-1$ telescopes the $\Delta$ terms:
\[
\Delta_T\le \Delta_0-\frac{\alpha}{2L}\sum_{t=0}^{T-1}\E\!\bigl[\norm{\grad f(x_t)}^{2}\bigr]
+\frac{L\alpha^{2}\sigma_g^{2}T}{2\mu^{2}}.
\tag{13}
\]
Rearranging,
\[
\frac{1}{T}\sum_{t=0}^{T-1}\E\!\bigl[\norm{\grad f(x_t)}^{2}\bigr]
\le\frac{2L}{\alpha T}\Delta_0
+\frac{L^{2}\alpha\sigma_g^{2}}{\mu^{2}}.
\tag{14}
\]
[Step 10]

\noindent\textbf{Step 11. Incorporate Hessian variance.}  

So far we have used only the gradient variance. To capture the effect of the stochastic Hessian we refine Lemma \ref{lem:hessian_inv}. Write $\widehat H_t = H(x_t)+\zeta_t$ where $\E[\zeta_t\mid x_t]=0$ and $\E\|\zeta_t\|_{F}^{2}\le\sigma_H^{2}$ (Assumption \ref{ass:hessian}).  
Using the matrix identity
\[
\widehat H_t^{-1}=H(x_t)^{-1}-H(x_t)^{-1}\zeta_t H(x_t)^{-1}+R_t,
\]
with remainder $\|R_t\|\le \frac{\|\zeta_t\|^{2}}{\mu^{3}}$ (derived from Neumann series, see e.g. \cite{matrix-perturbation}), we obtain
\[
\bigl\|\widehat H_t^{-1}\bigr\|
\le\frac{1}{\mu}+\frac{\|\zeta_t\|}{\mu^{2}}+ \frac{\|\zeta_t\|^{2}}{\mu^{3}}.
\tag{15}
\]
Taking expectation and using $\E\|\zeta_t\|^{2}\le\sigma_H^{2}$ yields
\[
\E\!\bigl[\bigl\|\widehat H_t^{-1}\bigr\|\bigr]
\le\frac{1}{\mu}+\frac{\sqrt{\sigma_H^{2}}}{\mu^{2}}+\frac{\sigma_H^{2}}{\mu^{3}}
\le\frac{2}{\mu}+ \frac{\sigma_H^{2}}{\mu^{3}}.
\tag{16}
\]
Plugging the bound (16) into (7) and (8) replaces $1/\mu^{2}$ by $C_H:=\bigl(\frac{2}{\mu}+ \frac{\sigma_H^{2}}{\mu^{3}}\bigr)^{2}$. Consequently the coefficient in front of $\norm{\grad f(x_t)}^{2}$ in (9) becomes
\[
\frac{\alpha}{L}-\frac{L\alpha^{2}C_H}{2}.
\tag{17}
\]
Repeating Steps 8–10 with $C_H$ yields the full bound
\[
\frac{1}{T}\sum_{t=0}^{T-1}\E\!\bigl[\norm{\grad f(x_t)}^{2}\bigr]
\le
\frac{2L}{\alpha\mu T}\bigl(f(x_0)-f^{\star}\bigr)
+\frac{2L\alpha\sigma_g^{2}}{\mu}
+\frac{2L^{2}\alpha^{2}C_H}{\mu^{2}}.
\tag{18}
\]
Since $C_H\le\bigl(\frac{2}{\mu}+\frac{\sigma_H^{2}}{\mu^{3}}\bigr)^{2}\le\frac{4}{\mu^{2}}+\frac{2\sigma_H^{2}}{\mu^{4}}$, the last term can be simplified to the expression stated in Theorem \ref{thm:main}.  
[Step 11]

\noindent\textbf{Step 12. Final simplification.}  

Bounding $C_H$ by $\frac{4}{\mu^{2}}$ and absorbing constants yields
\[
\frac{1}{T}\sum_{t=0}^{T-1}\E\!\bigl[\norm{\grad f(x_t)}^{2}\bigr]
\le
\frac{2\bigl(f(x_0)-f^{\star}\bigr)}{\alpha\mu T}
+\frac{2\alpha L\sigma_g^{2}}{\mu}
+\frac{2\alpha^{2}L^{2}\sigma_H^{2}}{\mu^{2}},
\]
which is exactly the claim of Theorem \ref{thm:main}. \hfill $\blacksquare$

\section*{Rate Derivation (With Constants)}
Set $\alpha = \Theta\!\bigl(T^{-1/2}\bigr)$, e.g. $\alpha = \frac{c}{\sqrt{T}}$ with $c\le\min\{1,\mu/L^{2}\}$.  
Plugging into the bound:
\[
\frac{1}{T}\sum_{t=0}^{T-1}\E\!\bigl[\norm{\grad f(x_t)}^{2}\bigr]
\le
\underbrace{\frac{2c^{-1}L\bigl(f(x_0)-f^{\star}\bigr)}{\mu}}_{\displaystyle O(T^{-1/2})}
\;+\;
\underbrace{\frac{2cL\sigma_g^{2}}{\mu}}_{\displaystyle O(T^{-1/2})}
\;+\;
\underbrace{\frac{2c^{2}L^{2}\sigma_H^{2}}{\mu^{2}}}_{\displaystyle O(T^{-1})}.
\]
Thus the dominant term decays as $O(T^{-1/2})$, matching the optimal rate for stochastic first‑order methods, while the Hessian‑variance term decays faster $O(T^{-1})$.

\section*{Special Cases}
\begin{itemize}
\item \textbf{Exact Hessian ($\sigma_H^{2}=0$).} The bound reduces to the classical stochastic Newton rate with only the gradient‑variance term.
\item \textbf{Large Hessian batch ($b_H\to n$).} By the law of large numbers $\sigma_H^{2}=O(1/b_H)$, so the third term becomes negligible for modest $b_H$.
\item \textbf{Convex $f$.} If $f$ is additionally $\mu$‑strongly convex, one can replace $\norm{\grad f(x_t)}^{2}\ge 2\mu\bigl(f(x_t)-f^{\star}\bigr)$ to obtain a linear convergence factor $1-\alpha\mu/L$ in expectation.
\end{itemize}

\begin{thebibliography}{9}
\bibitem{matrix-perturbation}
R. Bhatia, \emph{Matrix Analysis}, Springer, 1997.
\end{thebibliography}

\end{document}